% SPDX-FileCopyrightText: 2026 David Shirvanyants
% SPDX-License-Identifier: CC0-1.0

\chapter{Orientation Optimization}
\label{chap:optimize}

Optimize searches for a rotation which reduces the support required by a
model. Selected models are processed independently, in document-list order.
For every candidate orientation, a copy of the transformed triangle mesh is
rotated about its bounding-box center, sliced with the current slicing
settings, and passed through the Detect algorithm. The score $S$ is the sum
of unsupported polygon area over all layers; lower scores are better.

\section{Randomized Descent}

The unmodified orientation is scored first so that the user immediately sees
a baseline. The configured number of attempts are then distributed among the
worker threads. The first attempt starts at the baseline orientation. Each
additional attempt starts from a uniformly sampled random three-dimensional
rotation, which reduces dependence on any one local minimum.

Within an attempt, the optimizer chooses a random unit rotation axis and
applies a five-degree step. It accepts the candidate only when the score
decreases by more than the convergence tolerance. An accepted orientation
becomes the starting point for the next step and resets the failure count. A
rejected candidate is followed by another random axis. After twelve rejected
directions, the angular step is halved. Descent stops when the score is no
greater than the convergence tolerance, or when twelve directions also fail
at the minimum step of 0.5 degrees.

\begin{figure}[htbp]
    \centering
    \begin{tikzpicture}[
        node distance=3mm,
        every node/.append style={font=\small}
    ]
        \node[algstep] (start) {Start from the baseline or a random orientation};
        \node[algstep, below=of start] (perturb) {Rotate five degrees about a random axis};
        \node[algstep, below=of perturb] (score) {Slice, detect unsupported area, and calculate $S$};
        \node[algdecision, below=4mm of score] (better) {Improvement exceeds tolerance?};
        \node[algstep, below left=5mm and -5mm of better,
            text width=0.30\textwidth] (accept)
            {Accept candidate and publish a new global best};
        \node[algstep, below right=5mm and -5mm of better,
            text width=0.30\textwidth] (reject)
            {Count failure; halve the step after twelve failures};
        \node[algstep, below=14mm of better] (finish)
            {Finish at the stopping condition; begin another attempt if available};
        \draw[algarrow] (start) -- (perturb);
        \draw[algarrow] (perturb) -- (score);
        \draw[algarrow] (score) -- (better);
        \draw[algarrow] (better) -| node[pos=0.25, above] {yes} (accept);
        \draw[algarrow] (better) -| node[pos=0.25, above] {no} (reject);
        \draw[algarrow] (accept.south) -- (finish.north west);
        \draw[algarrow] (reject.south) -- (finish.north east);
    \end{tikzpicture}
    \caption{One orientation-search attempt. Repeated accepted and rejected
    steps form the local descent before the attempt finishes.}
\end{figure}

All workers share the best score and rotation under synchronization. When a
worker finds an improvement, the GUI applies that rotation on top of the
model's original transform for this run; source vertices are never rewritten.
The status bar reports completed attempts and the current best score. Stop and
model revisions are checked between expensive candidate evaluations, which
prevents a cancelled or stale result from being installed.

Random restarts make the search practical for a score surface with many local
minima, but the procedure is heuristic. It does not prove that the reported
orientation is globally optimal, and separate runs may produce different
results.
